\noindent Experimental evidence suggests that delegation can shift responsibility for adverse outcomes onto human intermediaries. We ask whether the same responsibility-avoidance motive applies to delegation to algorithms. In a preregistered online experiment with $322$ participants, a decision-maker chooses to make a payoff-relevant prediction for a matched recipient or delegate it to an algorithm, calibrated to perform equally well. We vary between subjects whether the recipient can impose punishment. The prospect of punishment reduces delegation by $18.3$ percentage points, from $60.8\%$ to $42.5\%$, reversing the direction that the responsibility-shifting account predicts. Delegation offers no protection in return. Recipients punish delegated and self-made predictions alike and condition punishment on the realized outcome. We discuss several candidate mechanisms for this reversal.